\documentclass[../main.tex]{subfiles}
\begin{document}
%==========================================TIÊU ĐỀ TẬP========================================
\begin{table}[h]
  \begin{adjustwidth}{-1cm}{}
    \begin{tabular}{>{\centering\arraybackslash}p{6cm}|>{\raggedright\arraybackslash}p{12.5cm}}
      \multirow{5}{6cm}
      % Cột bên trái: ÉPISODE và số tập
      {\\[-40pt]
        \flaregothic\fontsize{20pt}{20pt}\selectfont \centering
        \color{eptype}ÉPISODE\color{black}\\
        \burbank\fontsize{72pt}{72pt}\selectfont
        \color{epnum}127\color{black}\\[6pt]
        \cafeteria\fontsize{20pt}{20pt}\selectfont\color{epnum}(10-7)\color{black}
        \\[18pt]
        % \flaregothic\fontsize{14pt}{14pt}\selectfont
        % \color{epattr}BIENVENUE, \uline{\color{Banpire} Mel} !
        % \flaregothic\fontsize{18pt}{18pt}\selectfont
        % \color{epattr}ÉPISODE PERDU
        \color{black}
      }
       &
      % Dòng thứ nhất: tên tập tiếng Nhật
      {\fotskip\fontsize{18pt}{18pt}\selectfont もしもプロデューサーになったらどうなるのか？
      }            \\[6pt]
      % Dòng thứ hai: tên tập tiếng Anh
       & {\cafeteria\fontsize{18pt}{18pt}\selectfont 21st Century Sheep}         \\[4pt]
      % Dòng thứ ba: tên tập tiếng Việt
       & {\vnmsans\fontsize{18pt}{18pt}\selectfont Phi vụ 60 giây}               \\[8pt]
      % Dòng thứ tư: Nhân vật xuất hiện
       & {\sen{\fontsize{14pt}{14pt}\selectfont Track used: Chronos (とき)}}
      \\[6pt]
      % Dòng cuối cùng: Thời gian diễn ra của tập (thường sẽ là ngày phát hành)
       & {\continuum\fontsize{16pt}{16pt}\selectfont Ngày phát hành: 17/10/2021} \\[18pt]
    \end{tabular}
  \end{adjustwidth}
\end{table}
%=========================================TRUYỆN CHÍNH========================================
\color{black}

Hãy tưởng tượng rằng nếu các cô gái \hll\ chúng ta cùng nhau viết một kịch bản cho \textit{hologra}, mọi chuyện sẽ như thế nào? Một thế giới tràn ngập hỗn loạn, những ý tưởng bộc phát, và logic thì bị ném ra ngoài cửa sổ từ lâu. Và đúng thế, đó chính là chuyện của những ngày trước, và hôm nay\dots\ cũng không phải ngoại lệ.

Và\dots\ ờm\dots\ họ chỉ còn đúng sáu mươi giây để làm chuyện đó. Sáu mươi giây cho toàn bộ kịch bản.

``''\textbf{Khoan đã, cái này có trong kịch bản đâu!}'' một giọng hoảng hốt vang lên giữa trường quay, mà không ai khác chính là Watame, hôm nay bất đắc dĩ trở thành quản lý tạm quyền do Kuon bị ``tước ghế'' vì lý do không rõ ràng. ``Chúng ta chỉ có sáu mươi giây để làm thôi sao!?''

``Tôi xác nhận là đúng,'' một người bước ra từ hậu trường, tay cầm bản thảo lộn xộn toàn chữ là chữ. Đó là Cường, người vừa là tác giả, biên dịch viên, và giờ thì có vẻ đang kiêm luôn vai trò phản diện. ``Muốn đổ lỗi thì cứ đổ cho cái kịch bản này. Tôi chỉ là người ghi lại thôi, đừng bắn người đưa tin chứ.''

Kuon đứng khoanh tay, khuôn mặt lạnh tanh nhưng rõ ràng là đang kìm chế một cái nhíu mày. ``\vie{Nhưng nếu thế thì\dots\ chẳng phải là ngớ ngẩn quá hay sao? Một phút? Ông có biết một phút là bao nhiêu không!?}''

``\vie{Tôi biết,}'' Cường nhún vai đáp, ``\vie{nhưng mà tôi tưởng \textit{hologra} nào chả thế?}''

``\vie{\dots Tôi cũng chịu ông rồi đấy,}'' cậu Tổng - giờ thất nghiệp tạm thời - buông một tiếng thở dài, quay sang nhìn máy quay như thể đang tìm một câu trả lời từ vũ trụ. ``\vie{Chúng tôi từng làm mấy thứ còn vô lý hơn, nhưng ít ra thì vẫn được tập trước\dots}''

``\vie{Tôi cũng đang bó tay với chính bản thân mình đây,}'' tôi đáp lời, giọng đầy bất lực nhưng cũng có phần khoái trá.

Ở phía sau, Kanata đang đứng trên một cái ghế đẩu để với tới cục đạo cụ hình thiên thạch, thì thầm, ``Ờm\dots\ cho tớ hỏi nhỏ\dots\ cái đồng hồ này thực sự lơ lửng thật sao? Cái này nặng phết ấy.''

Lamy thì đang lúi húi thay trang phục lần thứ tư trong vòng mười phút. ``Cái bộ này là trang phục nhà thám hiểm hay là\dots\ phục vụ quán nhậu vậy? Sao lại có tạp dề ren!?''

``Cứ mặc bộ trang phục thường ngày là ổn rồi,'' Cường nói. ``Nội dung của tập này cũng không có gì quá đặc sắc đâu.''

Watame ôm đầu, rên rỉ: ``Không ai nghe tôi nói gì cả. Tôi là quản lý mà\dots\ ít nhất thì hôm nay là vậy\dots''

``Tốt rồi,'' cậu nói, lật sang trang cuối của bản nháp, ``chuẩn bị ghi hình đi. Sáu mươi giây bắt đầu\dots\ ngay bây giờ.''

\il

Bên trong một căn phòng khép kín ở tầng ba văn phòng, ánh đèn huỳnh quang trắng hắt xuống chiếc bàn gỗ lộn xộn tài liệu, dây cáp, và vô số lon nước tăng lực đã uống cạn. Căn phòng này, vốn được cải biên thành một studio thu nhỏ, đang vang lên tiếng bàn phím lách cách, tiếng chuột nhấp liên tục, và cả những tiếng than thở rời rạc không đầu không đuôi. Đã sát đến mười hai giờ đêm, nhưng tập phim \textit{hologra} mà họ dự định phát sóng vào chiều mai\dots\ vẫn chưa đâu vào đâu.

Bốn người đang làm việc cật lực: Kanata, đang đứng vắt vẻo trên một cái ghế quay, liên tục tua đi tua lại thước phim; Lamy, vừa gõ phụ đề vừa lầm bầm gì đó về Rushia và cỏ bạc hà; Kuon, người lẽ ra là quản lý chính thức nhưng hiện bị tước quyền vì ``tự ý thay đổi nhạc nền bằng organ điện tử''; và cuối cùng là Watame, người tạm thời ngồi vào ghế chỉ đạo với tấm bảng trắng ghi nguệch ngoạc dòng chữ: ``Tập phhim 60 giây - không gian căng thẳng - \textbf{CÒN LẠI MỘT PHÚT!}''

``Gay (go) rồi đấy\dots'' cậu con trai nhíu mày, giọng trầm xuống như thường lệ, ``Chúng ta không còn nhiều thời gian đâu.''

``Vậy chúng ta sẽ làm gì?'' nàng Yêu tinh tuyết hỏi, tay vẫn đang rê chuột để chọn hiệu ứng chuyển cảnh cho đoạn Rushia nhìn vào tủ lạnh rồi đóng cửa lại không lý do.

Trên màn hình, các đoạn phim hiện tại hầu hết chỉ là tư liệu hằng ngày, cảnh nàng phapsw ăn bánh, đi ngang qua hành lang, hoặc lẩm bẩm nói chuyện một mình trước gương. Không đoạn nào đủ kịch tính để mở đầu một tập phim trong vòng đúng một phút cả.

Nàng Cừu nheo mắt nhìn màn hình rồi đứng bật dậy: ``Như thế này không ổn đâu!''

``Sao vậy, trưởng phòng?'' nàng Thiên sứ nghiêng đầu hỏi, gãi gãi sau gáy.

Cô nàng Cừu - vẫn mặc bộ đồ cosplay kỳ quặc với một chiếc khố trắng quấn quanh người kiểu ``samurai thời kỳ Heian'' (không rõ lấy ở đâu ra) - khoanh tay lại, giọng nghiêm nghị không phù hợp với hình ảnh: ``Chúng ta phải xong trước giờ G!''

``Nhưng mà\dots\ ta chỉ còn lại một phút thôi!'' Kanata kêu lên, chỉ tay về chiếc đồng hồ treo tường đã mấp mé chỉ đến mười hai giờ đêm.

\img{10-7a}

``Anh nghĩ là ta không đủ thời gian đâu\dots'' cậu nói, giọng mệt mỏi. Cậu đã ngồi trước máy hơn mười bốn tiếng và tai vẫn còn ong ong vì tiếng chèn nhạc sai tông của chính mình.

Watame quay ngoắt lại, giậm chân một cái: ``\textbf{Tôi không cần biết, các người làm thế nào thì làm! Tôi là quản lý hôm nay, và tôi nói: phải ra được một phim!}''

Không còn cách nào khác, cả bốn người đành phải ngồi xuống, vắt óc suy nghĩ cho một kịch bản có thể\dots\ vừa quay, vừa dựng, vừa lồng tiếng trong sáu mươi giây. Không cốt truyện, không cảnh quay sẵn, chỉ có niềm tin mơ hồ rằng \textit{hologra} chưa bao giờ\dots\ cần logic.

Trên góc phải màn hình, đồng hồ bắt đầu đếm ngược {\digi 1:00} bắt đầu.

\ct

Thực ra thì họ chỉ còn lại 55 giây. Nhưng vì lý do ``kỹ thuật kịch bản'', ta cứ coi như đó là thời điểm họ bắt đầu làm việc (ghi chú từ Cường - vâng, là tôi đây).

Trong căn phòng chật hẹp giờ đã trở thành chiến trường sáng tạo, Lamy và Kuon vẫn đang cật lực xử lý video: tay chạy chuột, mắt đảo liên tục giữa các lớp phim, hiệu ứng, âm thanh. Cạnh đó, Kanata thì vừa lật giấy nháp vừa gõ gấp rút trên bàn phím, cố gắng vá víu một kịch bản mà chính cô cũng không chắc là mình đang viết gì. Còn ở phía sau, Watame - người đeo kính râm giữa nửa đêm như thể đang đi tuần tra hộp đêm - khoanh tay đứng giám sát tất cả, vẻ mặt nghiêm túc đến mức\dots\ khó tin được là một cô cừu.

\img{10-7b}

``Tiến trình thế nào rồi?'' nàng Cừu hỏi, giọng thấp nhưng đầy quyền uy.

``\textbf{Á!}'' Lamy giật bắn người, quay đầu lại như học sinh bị gọi lên bảng. ``C-Còn ba cảnh nữa!''

``Đang xử lý đoạn Koro-chan và Peko-chan tranh nhau lấy bóng,'' Kuon nói, tay vẫn nhấp chuột liên hồi như một lập trình viên đang chạy hạn chót qua đêm.

Watame gật đầu, đẩy kính râm lên một chút. ``Khẩu độ ống kính tốt nha.''

Không nói thêm một câu gì, cô bước lại chỗ Kanata vẫn đang lúi húi vừa viết vừa dựng, trông như sắp nhập viện vì stress. Watame cúi xuống, kéo nhẹ kính xuống sống mũi, nói một cách lạ kỳ ngọt ngào: ``Ấy? Gượm đã nào.''

Cô chỉ tay về phía màn hình, nơi hiện ra một cảnh quay Rushia đang ném một quả bóng chày ra xa với vẻ mặt tức giận. ``Thay quả bóng đó bằng con dao đi.''

Cả căn phòng chợt yên lặng.

Nàng Thiên sứ ngẩng phắt đầu, mắt trợn tròn: ``Ơ kìa!?''

Kuon thì gần như nghẹn cả câu: ``Không kịp đâu em ơi\dots\ Có gì dùng nấy thôi-''

``Tôi hỏi cô,'' nàng Cừu quay sang nhìn cô bạn cùng Thế hệ, giọng nghiêm như sĩ quan, ``không phải cậu. Còn cậu,'' cô nhìn Kuon như thể nhắc nhở anh về vị trí bị giáng chức tạm thời, ``tập trung chỉnh sửa phần Pekora và Korone cãi nhau đi.''

``E-Em sẽ thử!'' Kanata vội vàng đáp, mắt láo liên như đang tìm đường thoát thân.

``Đúng là KanKans của ta\~{}'' Watame nở nụ cười đắc ý, sau đó còn rướn người ra trêu Kanata bằng giọng trẻ con, ``Lần sau tôi sẽ bao cho cô\~{}''

\img{10-7c}

Trên màn hình lớn treo tường, đồng hồ chỉ về {\digi 0:42}. Hơn mười giây đã trôi qua.

Nàng Thiên sứ khẽ gật đầu, miệng thì ``cảm ơn'' nhưng mặt thì không giấu nổi vẻ tức tối, tay run run lướt chuột chỉnh các khung hình, chẳng hiểu là do mệt hay do ức chế.

``Cảnh này cũng phải sửa chút đi,'' Watame lại chỉ vào một phân cảnh khác. Cô nàng thiên sứ gần như há hốc.

``\textbf{Hở nữa!?}'' Kanata kêu lên.

Cấp trên của cô không trả lời trực tiếp, thay vào đó chỉ gật nhẹ rồi bắt đầu ``mách nước'' với tốc độ nói như rap.

``Các cô gái phải để kề nhau (juxtapose) rồi phóng to ra!''

``D-Dạ!''

``Cắt cảnh đột ngột sang Hà Nội!''

``Cái này anh vừa làm xong rồi,'' Kuon chen vào, vẫn mắt không rời màn hình.

``Chân thực hóa (vérité) cảnh trong cảnh, chồng các cảnh (superimpose) họ bị cảm lên nhau\dots''

``Chậm thôi, em ơi!'' cậu thở gấp.

``\dots dùng dolly shot\footnote{Kỹ thuật cảnh quay phóng nhỏ chủ thể khi tiến gần và phóng to khi ra xa chủ thể.} ở đây\dots''

Nàng Thiên sứ vẫn chỉnh sửa, nhưng trán đã bắt đầu đầm đìa mồ hôi, mặt cô nàng tái nhợt.

``\dots khuôn tô (gobo) cảnh tiền xu này, overcrank\footnote{Kỹ thuật tăng fps (số khung hình trên giây) để khi cảnh đó phát bình thường, người ta sẽ cảm tưởng nó được tua chậm lại.} cái hộp gối, phóng to cái phần giới thiệu ở cuối, rồi chốt lại bằng một câu: `Blessed you' nha!''

Tới mức này, Kanata không chịu nổi thêm nữa. Cô gần như hét lên, cắt luôn cả lời của Watame: ``\textbf{Thôi!!} Chị chẳng hiểu em đang nói cái gì cả!!''

``Thật,'' Kuon gật đầu tán đồng, ``đến anh chạy đồ án cũng chưa bao giờ gấp gáp đến thế này\dots''

Bỗng chốc, Watame chỉ tay về một khung cận cảnh của Rushia với ánh sáng lung linh, gương mặt mờ mờ, khiến cả ba người còn lại cùng dừng tay và nhìn.

\img{10-7d}

``Và rồi, ta có RuRu như thế này!'' nàng Cừu tạo dáng, giơ một tay lên cao như tượng Nữ thần Tự do. Miệng nở nụ cười ranh mãnh: ``Pơ-phệch.''

Lamy khẽ nở nụ cười cam chịu. Kanata thì rít qua kẽ răng: ``\hl{Em sẽ giết ả ta\dots}''

Kuon thì nói khẽ, gần như lẩm bẩm với màn hình: ``\hl{Thật\dots{} Sau vụ này chắc anh cũng xin bỏ việc trông xưởng luôn\dots}''

Bỗng nhiên, giữa lúc mọi người vẫn đang loay hoay với các khung hình và chỉnh sửa âm thanh, một giọng lo âu vang lên từ phía bàn dựng.

``Quản lý ơi\dots\ ta sắp hết thời gian rồi!'' nàng Yêu tinh tuyết, quay phắt đầu lại, mắt mở to như thể vừa phát hiện ra một quả bom hẹn giờ ngay trước mặt.

``\textbf{Cái gì chứ!?}'' Watame hét lên, hai chân như bật lò xo.

``Chính xác là còn {\digi 0:14} thôi,'' Kuon nói, giọng đều đều nhưng có phần lạnh gáy, nhưng không phải vì deadline, mà vì sợ phản ứng của cô cừu đang đứng sau lưng.

Không nói không rằng, Watame phóng như tên bắn về phía bức tường, nơi chiếc đồng hồ điện tử treo lơ lửng đang nhấp nháy đỏ rực từng giây một. Nó hiện lên con số {\digi 0:13}, rồi {\digi 0:12}\dots

``\textbf{Còn khuya nhá!!}'' cô hét lên, rồi lao tới\dots\ bám lấy chữ số 0 ở hàng giây như thể có thể ngăn cản bánh răng thời gian.

``\textbf{Giữ lấy con số chết chóc này!}'' Watame hét, tay bám chặt mép số, mắt long lanh như đang đóng một cảnh kịch Shakespeare giữa rạp phim hành động.

Cả ba người còn lại, Lamy, Kanata và Kuon, chỉ biết đứng im nhìn cảnh tượng hoang đường đó.

``\dots Ờm\dots'' cậu con trai chậm rãi lên tiếng, ``\dots anh nghĩ là em đã bị muộn rồi đấy.''

Và đúng lúc ấy, số 1 ở hàng chục giây từ từ biến mất, nhường chỗ cho 0. Watame quay đầu lại, đôi mắt cô trợn tròn nhìn từng pixel đổi màu như nhìn thấy thế giới sụp đổ.

\begin{figure}[H]
  \centering
  \begin{subfigure}[t]{0.45\textwidth}
    \centering
    \includegraphics[width=8cm]{../../../../Image/Book?/10-7e1.png}
    \caption{Watame tuyệt vọng bám lấy số 0 ở hàng đơn vị, như là nỗ lực cuối cùng để ``dừng thời gian'' lại\dots}
  \end{subfigure}
  \quad
  \begin{subfigure}[t]{0.45\textwidth}
    \centering
    \animategraphics[autoplay,loop,width=8cm]{12}{../../../../Image/Book?/10-7e2/10-7e2.}{1}{67}
    \caption{\dots chỉ để việc thời gian vẫn trôi, số 1 ở hàng chục giây dần hóa về 0 trước sự ngỡ ngàng của nàng Cừu.}
  \end{subfigure}
\end{figure}

``\textbf{Không!!!!}'', cô thốt lên trong nỗi tuyệt vọng pha chút bi hài, buông tay khỏi màn hình trong một tư thế sụp đổ.

\ct

Mười hai giờ hai mươi giây.

Trong khoảnh khắc mười giây cuối cùng vừa rồi, Kuon đã kịp lúc bấm nút gửi bản dựng cuối cùng lên hệ thống nội bộ.

Một tiếng *PING!* vang lên nhỏ xíu giữa căn phòng lặng ngắt, cùng lúc tiếng hét đau đớn của Watame vừa dứt.

``Anh đã nộp bản thảo rồi đấy\dots'' Kuon nói, tay rời khỏi chuột, ngả người ra sau như thể vừa hoàn thành một nhiệm vụ sinh tồn cấp độ khó.

``Còn gì nữa đâu mà khóc với sầu đây\dots'' Watame nói, giọng bây giờ đã chuyển sang lạc lõng như một nhân vật phụ vừa bị loại khỏi vai chính.

Cả bốn người ngồi im trong mấy giây, rồi đồng loạt quay ra màn hình, nơi thước phim vừa được nộp cho A-chan đang phát lại: cảnh Rushia tung quả bóng chày - không phải là con dao như lúc đầu do không còn thời gian -  vào khoảng không trong tiếng *VÚT* lồng nhạc sử thi.

Kanata và Lamy ngồi bất động, mặt cúi gằm như đang kiểm điểm trước hội đồng kỷ luật. Kuon thì ôm trán, thở dài một hơi dài không có hồi kết. Còn Watame ngồi xổm, miệng cười mỉm, như thể nỗi đau vừa rồi không là gì với cô: ``Vậy\dots\ tôi đã không làm gì sai\dots\ nhỉ?''

Lập tức, cả ba người còn lại quay sang cô, đồng thanh hét: ``\textbf{Quá sai luôn chứ còn gì nữa!!}''

\img{10-7f}

Nhưng cô nàng Cừu chẳng để tâm. Nhún vai, cô đứng dậy, phủi bụi không tồn tại trên váy mình, rồi bước chậm về phía cửa. Dừng lại ở bậc cửa, cô quay đầu, mắt ánh lên vẻ gì đó\dots\ kỳ lạ.

``Có thể lần này chúng ta vẫn chưa xong\dots'' cô nói, giọng như một nữ anh hùng trong phim chiến tranh.

``Hả?'' Kuon ngẩng đầu.

``\dots nhưng bây giờ,'' cô tiếp tục, mắt nheo lại, ``chúng ta lại có kịch bản \textit{hologra} tiếp theo đây.''

``\textbf{Gì chứ!?}'' Kanata và Lamy thốt lên gần như cùng lúc.

Cô nàng Cừu xoay hẳn người lại, tay chống hông, ánh mắt sắc lẹm: ``Tôi đang nói đến chúng ta, những ngôi sao của sân khấu này.''

Cả hai cô gái sửng sốt đến mức há hốc miệng, như thể vừa chứng kiến một vụ sự cố ngoài ý muốn. Còn Kuon\dots\ chỉ thở dài thật sâu, rất sâu.

``\textit{Mình muốn nghỉ việc\dots}'' cậu lẩm bẩm. ``\textit{Thà chạy deadline đồ án còn thoải mái hơn so với việc phải đối đầu với cái bà Cừu dở hơi điên khùng này\dots}''
%==========================================TRUYỆN PHỤ=========================================
\omk

\begin{center}
  \uline{Sáng ngày hôm đó\\Văn phòng.}
\end{center}

``Và đó\dots'' Watame chống tay lên bàn, mắt mơ màng như thể vừa bước ra khỏi một cuộc chiến tranh thế giới, ``\dots là những gì đã xảy ra. Tối qua. Trong vòng sáu mươi giây. Với một cái đồng hồ biết đếm ngược, một con dao thay cho bóng chày, và ba nạn nhân suýt nữa đốt máy dựng phim.''

Căn phòng lặng trong vài giây. Marine, người vừa nghe xong toàn bộ câu chuyện, đặt tay lên cằm, gật gù đầy hào hứng.

``Ra là vậy!'' cô reo lên, ``Thì ra \hll\ vận hành như thế đó!''

``Khoan đã,'' Kuon nói, ngẩng đầu khỏi tập giấy đang đọc dở, ``em tưởng chuyện đó là bình thường à?''

``Chứ không đúng hả?'' Senchou hỏi lại. ``Mỗi lần xem \textit{hologra}, em cứ tưởng công ty mình kiểu gì cũng đang vận hành như một xưởng hoạt hình thời chiến: ai muốn làm gì thì làm, đến giờ G thì nộp!''

Cậu thở dài. ``Anh dám cá là không đâu. Chẳng ai làm việc kiểu `nước đến chân mới nhảy' như thế cả. Nhất là trong ngành này.''

``Nhưng chẳng phải các số \textit{hologra} đều có cùng một mẫu số chung như thế sao?''

``Không có nghĩa là em vơ đũa cả nắm như thế,'' cậu đáp, giọng có phần bất lực. ``Nói như em thì mọi người ở đây đều\dots\ ăn cám lợn hết à? Với cả\dots''

Trước khi cậu nói tiếp, nàng Cừu đã cắt lời, mặt tái mét như nhớ lại cơn ác mộng đêm qua. ``Đừng\dots\ đừng có lôi em vào chuyện này chứ\dots''

``Thề,'' Kuon nói, dựa lưng vào ghế, ``anh mà là người làm việc theo kiểu đó, chắc bây giờ đang là quản lý của cái đội sản xuất đêm hôm ấy rồi đấy. Và cũng đếch có chuyện để cho mọi người dựng kịch bản không đầu không cuối, cốt truyện thì lủng củng khó hiểu thế này đâu.''

Marine nghiêng đầu. ``Ể\dots\ Nhưng em tưởng Wata-chan hợp với công việc đó mà\dots''

``Em ấy mà làm thật thì công ty có ngày sụp luôn chứ chẳng đùa.''

``Đêm qua nếu không vì kịch bản từ bên chúng ta thì có khi bị Kuon-senpai với Kanata đấm em rồi không chừng!'' Watame vội vàng tiếp lời, tay giơ lên thủ thế phòng ngự như thể đang ngồi giữa phiên tòa.

``Và nữa\dots'' Kuon hạ giọng, ``kể mấy chuyện như vậy, người ta sẽ tưởng công ty mình có tác phong làm việc kém đấy.''

``Thế à\dots'' Marine tròn mắt, gãi gãi má. ``Hmm\dots''

\img{10-7g}

Rồi như thể một tia sáng lóe lên trong đầu, Senchou đột ngột chồm người ra phía trước, mắt long lanh: ``Vậy nếu em dùng thân thể này-''

``Không được,'' hai người còn lại đồng thanh cắt ngang, giọng cứng rắn như một cặp bố mẹ vừa bắt gặp con mình định nhảy vào hồ nước bẩn.

Marine phồng má, nhăn mặt nói nhịu: ``Em tưởng đó là ý hay mừ\dots''

Sau một lúc suy nghĩ, cô lại chuyển đề tài:
``À mà, em nghe nói ở mấy công ty sáng tạo, bộ phận nhân sự hay tuyển người theo kiểu ngẫu hứng ấy. Kiểu chọn vì `năng lượng tiềm tàng' chứ không cần CV hẳn hoi\dots''

Kuon liếc mắt. ``Ý em là kiểu họ\dots\ lấy tài năng ngẫu nhiên ấy hả? Cái đó thì phải hỏi ông CEO của chúng ta - chú YAGOO-san. Không thuộc chuyên môn của anh.''

Watame thì gật gù, lặng lẽ chăm chú nghe hai người nói, như thể đang cố nhớ lại thêm vài chi tiết để bổ sung cho câu chuyện kỳ quái của mình, dù ai cũng hiểu rằng, đáng lẽ nó không nên xảy ra.

Ngày hôm đó, không có chuyện gì bất thường cả.

Chỉ có Marine, Watame, và Kuon ngồi thêu dệt - hay đúng hơn, đính chính - những câu chuyện tưởng chừng như không thể xảy ra bên trong công ty \hll.

Không có gì đáng ngờ cả, YAGOO xin thề.

\ct

Quay trở về trường quay, nơi mà tác giả của bộ sách này đang trú ngụ, nơi treo đầy bản in cũ của các tập truyện, hồ sơ nhân vật, và một vài bức tranh minh họa vẽ dở dang, Cường đang ngồi xem lại hộp thư nội bộ thì cánh cửa bật mở.

``Cường-san!'' một giọng gọi vui vẻ vang lên.

Cậu ngẩng đầu lên. Là Watame, trong tay cô là một xấp giấy được in vội, mép giấy còn cong cong vì vừa được xé ra khỏi máy in. Theo sau cô là Lamy và Kanata, mỗi người đều mang vẻ mặt khó đoán: một người có vẻ tò mò, người còn lại thì hơi lo lắng.

``Đây!'' Wnàng Cừu đặt xấp kịch bản lên bàn, gõ nhẹ một cái như thể đang đặt một bản hiệp định lịch sử, ``Tập \textit{hologra} Rushia hôm nay. Bản hoàn chỉnh!''

Cậu cầm lấy tập giấy, lật vài trang đầu tiên. Rất nhiều đoạn viết tay, ghi chú bên lề bằng bút đỏ, khoanh tròn, mũi tên chỉ lung tung, và không thiếu những chỗ bị gạch bỏ rồi viết lại bằng ký hiệu chẳng giống ai.

``Để em giải thích nha,'' Watame nói, gần như không cần chờ Cường lên tiếng. ``Ở đoạn đầu, ta nên mở bằng một cú dolly shot xuyên qua hành lang, rồi cảnh Rushia hiện ra trong một khung cận mặt, ánh sáng từ bên trái nha, đừng từ trên xuống! Rồi đoạn con dao- à, em nhầm\dots\ quả bóng, phải có hiệu ứng motion blur 7 độ, để nhấn mạnh tốc độ! Và\dots''

Cô nàng thao thao bất tuyệt. Từ nhịp cắt dựng đến chuyển cảnh, từ độ bão hòa của cảnh bếp cho đến cách ``fade-out'' đổ bóng hình trái tim khi Rushia nhắm mắt. Cậu ngồi đó, vừa nghe vừa lật tiếp vài trang, mắt đảo nhẹ theo từng dòng chữ loằng ngoằng, đầu gật gù theo phép lịch sự.

Nàng Thiên sứ và nàng Yêu tinh tuyết đứng bên cạnh, ánh mắt hết liếc sang cậu tới sang nàng Cừu, rồi lại về cậu. Trong thâm tâm, cậu nghĩ thầm: ``\vie{\textit{Rõ ràng là họ đã nghe bài này ít nhất là ba lần rồi\dots}}''

``Rồi sau đó, ở đoạn kết,'' Watame tiếp tục, ``mình sẽ dùng một cú overcrank, để slow motion (tua chậm) ánh mắt cô ấy, và rồi lồng vào câu `Blessed you!' như em đã nói lúc nãy, nhớ không?''

Cường đặt tập giấy xuống bàn, thở dài lấy một hơi, ngẩng đầu nhìn cô: ``\dots Vậy chốt lại là em muốn chuyển thể nó vào truyện đúng không?''

Nàng Cừu chớp mắt: ``Vâng! Em nghĩ nó sẽ thành một chương ngoại truyện rất hay cho `Bài thơ lục bát về một con chó đang ngủ' ấy mà-''

Cậu gật đầu, nhưng rồi thở ra nhẹ một cái, chậm rãi nói: ``Vậy thì\dots\ mấy cái hiệu ứng hình ảnh em nói ấy\dots\ có dùng được đâu. Đây là \emph{tác phẩm văn học} mà.''

Câu nói vừa dứt, căn phòng lặng trong một giây.

``\textbf{Cái gì!?}'' cả Lamy lẫn Kanata đồng thanh quát, quay sang nhìn Watame bằng ánh mắt như sắp lôi cô đi họp kiểm điểm.

``Đó! Em đã bảo mà!'' nàng Thiên sứ rít lên. ``Ngay từ đầu em đã nói là mấy cái `motion blur 7 độ' của bà không áp được vào trang giấy đâu!''

``Watame-senpai\dots\ thật là hết thuốc chữa mà!'' nàng Yêu tinh tuyết thở dài ngán ngẩm. ``Chuyển thể truyện mà làm như dựng MV\dots''

Nàng Cừu cúi đầu, vẻ hơi quê, nhưng vẫn nói nhỏ: ``Nhưng\dots\ tôi nghĩ nó sẽ hay mà\dots''

Cường cầm lại bản kịch bản làm vội đó, ngả người vào ghế. Dù nét chữ hơi rối, vài đoạn gần như không thể đọc được, nhưng\dots\ nó vẫn có một năng lượng hỗn loạn rất riêng, một thứ năng lượng của \textit{hologra}.

Cậu lật sang trang thứ ba, rồi kẽ cười: ``\dots Nói thế thôi. Có khi anh sẽ nhờ GPT thử dựng lại nó, có thể sẽ giữ lại mấy cái chi tiết 7 độ với mấy cái điên rồ này á.''

Cả ba cô gái đều tròn mắt nhìn tôi.

``Nhưng,'' tôi nói thêm, ``lần sau đừng bắt nhân vật ném dao nữa. Nhất là Rushia.''

Watame ngập ngừng gật đầu. ``Dạ\dots\ Em sẽ thử\dots\ ném thứ khác\dots''

``\textbf{Không là không!!}'' Lamy và Kanata hét lên, một lần nữa. Cậu biên kịch phì cười với phản ứng có phần quyết liệt đó.

Đến cuối cùng, dù là ai dựng nên kịch bản \textit{hologra}, thì kết cục sẽ luôn chỉ có một mà thôi: một câu chuyện đầy tính ngẫu nhiên và hỗn loạn!

\end{document}